%% ============================================================================
%% VI. Experimental Plan and Preliminary Results              target: ~1.3 pages
%% Stage 1: experiments may be missing or preliminary, but the plan must show
%% how each hypothesis will be supported or falsified (reviewers score the
%% "feasibility of conclusion"). Stage 2: rename to "Experiments" and report
%% the experiments required by the Extension Charter.
%% Anonymity: no lab names or logos in photos or text.
%% ============================================================================
\section{Experimental Plan and Preliminary Results}
\label{sec:experiments}

\begin{figure}[t]
  \centering
  \figplaceholder{\linewidth}{4cm}{%
    \textbf{Fig.~4 (setup).} Robot arm with wrist force/torque sensor and
    gripper holding the Nitinol rod; sensorized ring (edge); RGB-D camera used
    for ground truth only.}
  \caption{Experimental setup.}
  \label{fig:setup}
\end{figure}

\subsection{Setup}
\label{sec:exp:setup}
\begin{outline}
  \item Yaskawa Motoman SDA10F (one arm), Bota SenseONE EtherCAT force/torque
    sensor at the wrist, parallel gripper \todo{model}, Nitinol rod
    \todo{length, diameter}, conductive-TPU sensorized ring, Intel RealSense
    D435 RGB-D camera for ground truth only \todo{tip marker / tracking method}.
  \item Planar (2D) experiments first, then spatial (3D).
\end{outline}

\subsection{Data collection and protocol}
\label{sec:exp:protocol}
\begin{outline}
  \item Random and structured gripper trajectories that cover contact and
    no-contact states over a range of deformations; each sample is
    $(\gpose,\wrench,\cpt,\tpose)$.
  \item Train/test split by trajectory (testing on unseen configurations);
    repeated trials to quantify sensor noise and hysteresis.
  \item Mitigations for sensor noise, inaccurate contact information, and
    synchronization errors: force/torque tare and calibration, time
    synchronization, camera--robot extrinsic calibration, repeated controlled
    trials.
  \item \todo{Dataset size, number of rods and edges, range of contact
    locations.}
\end{outline}

\subsection{Metrics and baselines}
\label{sec:exp:metrics}
\begin{outline}
  \item Estimation: tip position error [mm] and direction error [deg] on
    held-out data.
  \item Manipulation: success rate within $(\epsilon_p,\epsilon_\theta)$, final
    error, time to target, peak contact force.
  \item Baselines and ablations: inputs ($\gpose$; $\gpose{+}\wrench$;
    $\gpose{+}\wrench{+}\cpt$); model comparison; a rod-model (physics)
    estimator; contact-free and open-loop steering.
\end{outline}

\begin{table}[t]
  \centering
  \caption{How each hypothesis is tested}
  \label{tab:plan}
  \footnotesize
  \begin{tabular}{@{}p{0.06\linewidth}p{0.40\linewidth}p{0.44\linewidth}@{}}
    \toprule
     & Experiment & Supported if / refuted if \\
    \midrule
    \ref{hyp:reach}    & Reachable tip directions with vs.\ without contact
                         (simulation and hardware)
                       & Larger range with contact / no gain \\
    \ref{hyp:estimate} & Estimation error on unseen data; input ablation
                       & Error within tolerance and significant gain from
                         contact input / not \\
    \ref{hyp:steer}    & Steering to targets behind the edge vs.\ baselines
                       & Higher success, lower final error / not \\
    \bottomrule
  \end{tabular}
\end{table}

\subsection{Preliminary results}
\label{sec:exp:prelim}

\begin{table}[t]
  \centering
  \caption{Preliminary tip-pose estimation error (planar, held-out data)}
  \label{tab:prelim}
  \footnotesize
  \begin{tabular}{@{}lcc@{}}
    \toprule
    Inputs & Position [mm] & Direction [deg] \\
    \midrule
    $\gpose$                    & -- & -- \\
    $\gpose+\wrench$            & -- & -- \\
    $\gpose+\wrench+\cpt$       & -- & -- \\
    \bottomrule
  \end{tabular}
\end{table}

\begin{outline}
  \item \todo{First planar dataset: size, error statistics (\cref{tab:prelim}),
    and what they suggest about \ref{hyp:estimate}.}
  \item \todo{Reachability in simulation for \ref{hyp:reach}.}
\end{outline}

\subsection{Planned experiments (Stage 2)}
\label{sec:exp:planned}
\begin{outline}
  \item Spatial (3D) experiments; closed-loop steering trials
    (\ref{hyp:steer}); generalization to other rods and edge geometries.
\end{outline}
